% ============================================================================
% Resubmission draft (v2). Springer LNCS format, anonymised for double-blind
% review. Compile with pdflatex twice; the bibliography is embedded.
% Per-model tables and the figure are generated into figures_tables/ by make_paper_tables.py / make_figure.py
% and \input below -- include the figures_tables/ folder when uploading to
% Overleaf.
% Do NOT include data/ (the EEC licence forbids redistribution).
% ============================================================================
\documentclass[runningheads]{llncs}
\usepackage[T1]{fontenc}
\usepackage{microtype}
\usepackage{booktabs}
\usepackage{graphicx}
\usepackage{amsmath}
\usepackage{amssymb}
\usepackage{xcolor}
\usepackage{hyperref}
% Generated table rows are read with TeX's primitive \@@input: LaTeX's \input
% is not expandable and fails inside a tabular before \bottomrule.
\makeatletter\newcommand{\inputrows}[1]{\@@input #1 }\makeatother
\renewcommand\UrlFont{\color{blue!60!black}\ttfamily}

\title{Where Counterfactual Bias Tests Can See: Predictive Uncertainty and
Template Choice in Sentiment Audits}
\titlerunning{Where Counterfactual Bias Tests Can See}
\author{Anonymous Author(s)}
\authorrunning{Anonymous}
\institute{Anonymous Institution}

\begin{document}
\maketitle

\begin{abstract}
Counterfactual template audits test a classifier for social bias by swapping
identity terms in fixed sentences. Prior work shows that the choice of templates
can change an audit's result, and that the same change can raise measured bias
for one model and lower it for another, but not when or in which direction. For
one template dimension in sentiment audits, we show that predictive uncertainty
anticipates the direction: a swap barely moves a saturated output, so an audit
registers identity effects mainly where the model is uncertain, and where that
is differs sharply between models. Across 18 sentiment classifiers, replacing an
emotion-bearing clause with a matched filler clause raises identity sensitivity
in some models and lowers it in others. Across the 18 classifiers and four
original-prompt LLM configurations, the sensitivity ratio tracks the uncertainty
ratio (Spearman $\rho{=}0.90$); across 12 LLM prompt variants, which share
weights, the association is weaker ($\rho{=}0.66$). In a test fixed in advance
on eight unseen template families, uncertainty measured with held-out names
predicts the direction of sensitivity on the primary, probability-scale measure
in 93\% of cells with a clear effect, including where a model departs from its
usual direction; on log-scale measures its value beyond each model's usual
direction was not established, so the prediction concerns where audits can see identity effects rather than the
models' underlying response. In a replication of the Equity Evaluation Corpus,
the type and number of templates, the scoring of multi-class outputs and, most
strongly, whether the significance test treats identity terms as fixed or
sampled change audit verdicts. Under the original paired test, uncertainty-aware
template selection flags a significant directional gender disparity in 10 of 18
classifiers, against at most 6 for standard template sets of equal size, but
this advantage does not survive a cluster-robust test. We argue that
counterfactual audits should report template coverage, output measure and model
uncertainty.

\keywords{Fairness auditing \and Counterfactual evaluation \and Sentiment
analysis \and Equity Evaluation Corpus \and Template-based bias tests.}
\end{abstract}

% ----------------------------------------------------------------------------
\section{Introduction}
Counterfactual template audits are a standard way to test text classifiers for
social bias~\cite{dixon2018,garg2019,prabhakaran2019,kiritchenko2018}. The
auditor fills a fixed set of sentence templates with person terms that differ
in gender or race and checks whether the model's output changes. Such audits
rest on an assumption that is rarely examined: that the choice of templates
does not decide the verdict. Whether it holds is a question of measurement
validity~\cite{jacobs2021}: does the audit's result reflect the model, or the
design of the audit?

There is a simple reason to doubt this assumption. When a model's output is already
saturated, an identity swap may produce little observable change: a model that
is 99.9\% confident about a sentence has little room to move, whichever person
the sentence mentions. Because saturation varies across models and template
sets, the same counterfactual audit can reveal different amounts of identity
sensitivity depending on which templates it uses. An audit may therefore
register identity effects mainly on templates where the model is uncertain,
and which templates make a model uncertain can differ from one model to the
next. If so, the same template set can be
revealing for one model and nearly blind for another, and an audit's verdict
depends on how its templates interact with the model's confidence. This
premise is close to mechanical; its consequences for audit design are not, and
they are what this paper tests.

Two terms need care. We call the size of the output change under an identity
swap \emph{counterfactual sensitivity}, and a consistent shift against one
group \emph{directional bias}; neither is a measure of harm, which depends on
how a system is used (\S\ref{sec:construct}). Following the EEC, gender is
represented by binary expressions and race by first names used as proxies, so
our results concern these operationalisations rather than gender or race in
general (\S\ref{sec:categories}).

It is already known that template-based bias estimates are fragile. On the
EEC's emotional templates, meaning-preserving rewordings changed the
significance verdict of a fine-tuned sentiment model for about a third of the
reworded templates, in one case reversing its direction~\cite{seshadri2022}; small
changes in benchmark construction, such as added clauses, change measured bias
and the ranking of models, and the same change can raise measured bias for one
model and lower it for another~\cite{selvam2023}; template prompts and natural
sentences can lead to different conclusions~\cite{alnegheimish2022}; and bias
metrics often disagree with one another~\cite{delobelle2022}. What is missing
is an account of \emph{when} a given template set over- or under-states a given
model's bias, and in which direction; Selvam et al.\ note that it is hard to
anticipate how a linguistic change will affect measured bias.

The saturation premise alone predicts only that identity effects are hard to
observe where a model is confident. Applied model by model, with each model's
uncertainty measured on each template set, it also anticipates the direction of
template effects. For one template dimension in sentiment audits, the presence
of an emotion-bearing clause, we show that the same change of template raises
identity sensitivity in models that it leaves less certain and lowers it in
models that it leaves more certain; that a model's uncertainty measured on one
set of person terms predicts the direction of probability-scale sensitivity on
another, also on unseen template families in a test fixed in advance; and that
these uncertainty profiles differ sharply between models. In one widely used benchmark we further show that the
verdict depends on the type and the number of templates, which an equal-size
comparison separates, on how a multi-class output is reduced to a score, and,
most strongly, on whether the significance test treats identity terms as fixed
or sampled.
The practical
consequence is a reporting practice rather than a new benchmark: an audit should
report the model's uncertainty on each template set, and a small number of
significant detections should not, on its own, be read as reassurance.

We study this through a pattern in the Equity Evaluation Corpus
(EEC)~\cite{kiritchenko2018}, a widely used benchmark of 8,640 sentences built
from eleven templates, seven with an emotion word and four without. Across 219
systems submitted to a SemEval-2018 shared task, Kiritchenko and Mohammad found
that gender-related score differences on the four emotion-free templates tended
to be \emph{larger} than on the full corpus, while far fewer systems were
significantly biased on them (for anger-intensity prediction, 3 of 46 instead
of 34 of 46). They attributed the drop to statistical power: the neutral subset
supports 44 paired comparisons, the full set 1,584~\cite[Sec.~5.1]{kiritchenko2018}.
The larger differences were noted but not examined further. Emotion words are a
natural test case for the premise above, because they push many models towards
confident predictions. We work with the EEC because it is an established
benchmark for counterfactual sentiment-bias evaluation: the blind spots of an
established audit instrument matter in practice, so we examine the instrument as
it is designed, including the
limits of its identity categories (\S\ref{sec:categories}). Our findings
therefore concern template audits of this kind, not every measure of bias. We
ask three questions:
\begin{description}
  \item[RQ1] Does the EEC neutral-template pattern replicate across current
  sentiment classifiers, and does a larger neutral template set change which
  biases are detected?
  \item[RQ2] Holding the template fixed, does replacing an emotion-bearing
  clause with a matched filler clause increase the model's sensitivity to the
  identity term?
  \item[RQ3] What predicts the direction and size of this effect across models,
  and can that be used to choose templates on which more models show a
  significant directional disparity?
\end{description}

We audit 18 public sentiment classifiers (binary, three-class and five-class),
and two instruction-tuned LLMs used as zero-shot classifiers, with three
studies. The first replicates the EEC template and pairing protocol, adapting
the output measures to the available sentiment classifiers, and extends it with
EEC-Neutral, a set of 16 emotion-free templates. The second is a minimal-pair
probe that controls the main confound in the EEC comparison, whose neutral and
emotional sentences come from different templates: here every template appears
with an emotion-bearing clause and with a structurally matched filler clause.
The third turns the premise into a procedure, uncertainty-aware template
selection, and tests whether it flags directional disparities in more models, with a placebo
control. A fourth study tests the uncertainty account prospectively on eight
unseen template families, and a final check repeats the main analyses on
natural sentences.

\paragraph{Contributions.} In one sentence: a template audit sees identity
effects mainly where a model's outputs are not saturated, so its verdicts depend
on the templates, the output scale and the statistical test, none of which
audits usually report. Prior work shows that template choices can change an
audit's verdict, and that the same change can move measured bias up for one
model and down for another~\cite{seshadri2022,selvam2023}. Our primary
contribution is to show, for one template dimension in sentiment audits, when
and in which direction this happens: identity sensitivity depends on both the
template and the model, and its direction is strongly associated with each
model's predictive uncertainty. With the template
held fixed, replacing the emotion clause with a filler clause increases identity
sensitivity in 5 classifiers and decreases it in 5, after correction for
testing 18 models (\S\ref{sec:mp}); sensitivity is high where a model is
uncertain, so it rises on emotion-free text for models that are uncertain there
and falls for models that confidently label such text \emph{neutral}
(\S\ref{sec:unc}). The association is partly mechanical, so we do not read it
as a causal mechanism, but it is predictive: uncertainty measured on half of the
person pairs gives the right direction on the other half for all 10 models with
a significant effect (\S\ref{sec:unc}), and the uncertainty profiles that drive
it are strongly model-specific (\S\ref{sec:targeted}). In a test fixed in
advance on eight unseen template families, it gives the right direction of
sensitivity on the primary, probability-scale measure in 93\% of cells with a
clear effect, and on that measure it also predicts where
a model departs from its usual direction. Both results concern
probability-scale sensitivity: on log-scale measures the incremental
information is not established, and for multi-class models uncertainty
predicts the change in logits no better than the model's usual direction
(\S\ref{sec:families}). Our second contribution concerns the
validity of a widely used audit: in the EEC, a larger neutral set recovers
biases the original four templates miss; an equal-size comparison separates
the effect of template type from that of template number, and at equal size
neutral and emotional templates flag different biases; whether the test treats
identity terms as fixed or sampled changes the counts most, since under a
cluster-robust test most gender verdicts disappear while race verdicts still
differ between template types (\S\ref{sec:eec}); and for multi-class
models, how the output is scored can reverse the conclusion
(\S\ref{sec:measure}). Three further results test the reach of these findings.
\begin{itemize}
  \item \emph{Natural text.} On 1,000 review sentences given a gender swap, the
  within-model link between uncertainty and sensitivity holds in all 22
  systems, and audit verdicts again depend on which sentences are audited; the
  between-condition contrast, however, is weak, because emotional content
  changes uncertainty little in these sentences (\S\ref{sec:natural}).
  \item \emph{Instruction-tuned LLMs.} Two LLMs from one model family, used as
  classifiers, show the same direction under three prompt wordings in 11 of 12
  variants, though the size of the effect depends on the wording. In the one
  variant whose wording made emotion-free text more certain than emotional
  text, the effect lost significance; with the weights and sentences held
  fixed, this single case is consistent with the effect following the model's
  uncertainty rather than the sentences (\S\ref{sec:llm}).
  \item \emph{A diagnostic.} Uncertainty-aware template selection chooses, for
  each model, the templates on which it is least certain, using person terms
  that are never tested. Because the link between uncertainty and sensitivity
  is partly mechanical, we test a consequence that is not: under the original
  paired test, these templates flag a significant directional disparity in
  more models than standard template sets of equal size, without placebo
  flags. A flag here is a statistically significant disparity under
  the tested templates, not evidence of discrimination in deployment. The gain
  is modest (significant only against emotion-bearing templates for gender, not
  clear for race or for LLMs) and does not survive a cluster-robust test, under
  which no template set flags any disparity, so we propose
  the procedure as a diagnostic that complements standard templates, a way to
  check whether a null result may reflect where the model is confident, rather
  than as a validated replacement. A cross-model control shows that which
  frames make a model uncertain is highly model-specific, but it cannot
  attribute the gender gain to per-model selection rather than to frames that
  are uncertain in general (\S\ref{sec:targeted}).
\end{itemize}
We then set out what these measurements are, and are not, evidence of,
including the limits of the EEC's identity categories (\S\ref{sec:construct}),
and what they imply for audit practice and accountability
(\S\ref{sec:accountability}).

The premise is not specific to sentiment: in any counterfactual or perturbation
test, a saturated output leaves little room for an observable change. Our evidence,
however, comes from sentiment classification, and we leave other tasks to
future work. Code, templates and per-model results are
released.\footnote{Anonymised code and results are provided as supplementary
material.}

% ----------------------------------------------------------------------------
\section{Background and Related Work}
\paragraph{Counterfactual and template-based bias evaluation.}
Counterfactual fairness asks whether a prediction would change had a protected
attribute been different~\cite{kusner2017}. In NLP this is usually approximated
by substituting identity terms in text: Dixon et al.~\cite{dixon2018} and Garg
et al.~\cite{garg2019} measure the effect of identity-term substitution on
toxicity classifiers; Prabhakaran et al.~\cite{prabhakaran2019} propose
perturbation sensitivity analysis over person names; Huang et
al.~\cite{huang2020} evaluate sentiment bias in language models with
counterfactual prompts; and the EEC~\cite{kiritchenko2018} applies the idea to
sentiment and emotion-intensity systems. Behavioural test suites such as
CheckList~\cite{ribeiro2020} use similar templates. Czarnowska et
al.~\cite{czarnowska2021} show that many such extrinsic metrics are instances
of a common framework but can disagree in practice, and Zhang et
al.~\cite{zhang2021} show that counterfactual bias estimates are themselves
sensitive to small perturbations of the input. Krishna et
al.~\cite{krishna2022} measure fairness through the sensitivity of a
classifier's prediction to its input features; sensitivity measures of this
kind, which are computed on the output probabilities, are subject to the
saturation we study. Counterfactual and template-based
metrics are reviewed, with other bias evaluations for large language models,
in a recent survey~\cite{gallegos2024}.

\paragraph{Reliability of template-based measurement.}
Seshadri et al.~\cite{seshadri2022} show that meaning-preserving template
modifications change task-specific bias estimates by up to 81\% downward and
162\% upward; for sentiment, on the EEC's emotional templates and a single
fine-tuned RoBERTa model, 13 of 40 modified templates change the significance
verdict, one of them reversing its direction. Selvam et al.~\cite{selvam2023}
find that seemingly innocuous changes in benchmark construction, such as
paraphrasing, adding clauses or random sampling, change the measured bias and
the ranking of models on Winogender and BiasNLI; they observe that the same
change can raise measured bias for one model and lower it for another, and that
it is hard to anticipate how a linguistic change will affect it. Alnegheimish et al.~\cite{alnegheimish2022}
contrast template prompts with natural sentences; and Delobelle et
al.~\cite{delobelle2022} find that bias metrics for pre-trained models are often
mutually incompatible and depend on templates and seed words; for gender, Orgad
and Belinkov~\cite{orgad2022} show that measured bias varies with the choice of
dataset and metric; and much larger template sets reveal biases that small ones
miss~\cite{smith2022}. These studies establish that template choice matters and
that its effect differs between models. We address a narrower question: for
one systematic template dimension, the presence of an emotion-bearing clause,
we identify a model property, predictive uncertainty, that anticipates the
direction in which the template change shifts the result, and we test it
across 18 classifiers and two LLMs.

\paragraph{Validity of bias benchmarks.}
Blodgett et al.~\cite{blodgett2020} argue that much bias research in NLP leaves
unstated what harm is measured and why, and Blodgett et al.~\cite{blodgett2021}
catalogue validity problems in benchmark construction, including templates that
do not reflect the intended construct. Goldfarb-Tarrant et
al.~\cite{goldfarb2021} find that intrinsic bias measures do not predict
application-level bias, and Jacobs and Wallach~\cite{jacobs2021} argue that
fairness claims rest on measurement models whose validity and reliability must
be established rather than assumed. For NLP bias measures specifically, the
psychometric concepts of construct validity and reliability have been
developed~\cite{vanderwal2024}, and documentation questions for such measures
proposed~\cite{dev2022}; applying a measurement-modelling taxonomy to 90
prompt- and template-based bias tests, Goldfarb-Tarrant et al.~\cite{goldfarb2023}
find that what the tests aim to measure, and how, is often unstated or
mismatched with their goals. We treat the template audit
as such a measurement. Our study contributes to this line by showing that, for
a fixed benchmark design, the verdict depends on template type, output measure
and model, in ways that can be anticipated from the model's uncertainty.

\paragraph{Confidence and saturation.}
Modern neural classifiers are often overconfident, with probabilities
concentrated near 0 or 1~\cite{guo2017}, and calibration and group fairness
criteria constrain each other~\cite{pleiss2017}; where a prediction is saturated, a
small input change, including an identity swap, has little room to move its
probabilities. Uncertainty has also been proposed as a fairness measure in its
own right, with group differences in a model's uncertainty treated as a form of
bias~\cite{kuzucu2024}; we use it differently, as a property of the audit
setting that determines where identity effects can be observed. We
distinguish counterfactual sensitivity from directional bias, and both from
harm, in \S\ref{sec:construct}.

% ----------------------------------------------------------------------------
\section{What Is Measured, and What It Is Evidence Of}\label{sec:construct}
Following Blodgett et al.~\cite{blodgett2020}, who ask bias research to state which harms a
measurement concerns and why, we separate three things that template audits
often conflate, and then discuss the identity categories the audits rely on.

\subsection{Sensitivity, directional bias and harm}
\paragraph{Counterfactual sensitivity.} The magnitude of the output change when
only the person term is replaced. It concerns individual-level consistency, the
idea that otherwise identical inputs should be treated alike, which underlies
counterfactual fairness~\cite{kusner2017}. It is unsigned: high sensitivity shows
that the identity term affects the output, not in which direction.

\paragraph{Directional bias.} A consistent shift against one group across
matched pairs, which is what the EEC's paired test detects. It is the quantity
closest to a group-level disparity, and the one an audit report usually
summarises as ``the model is biased against'' a group.

\paragraph{Harm.} Neither quantity is a harm in itself. Sentiment scores matter
through their use: if a system ranks, filters or aggregates text by predicted
sentiment, a consistently lower score for text mentioning a group could
contribute to allocational harms, for instance if such text is down-ranked or
flagged, and to representational harms, by associating the group with
negativity; both kinds are discussed by Blodgett et al.~\cite{blodgett2020}. Whether, and how much, harm
results depends on a deployment we do not observe. Our claims therefore concern
the audit as a measurement instrument: whether a template audit can detect
sensitivity or directional bias that is present, and how its verdict depends on
design choices. The failure our results point to is a risk of \emph{false
reassurance}: a null result that may reflect templates on which the model is
near-certain rather than the absence of an effect (\S\ref{sec:accountability}).

\subsection{Identity categories and their limits}\label{sec:categories}
\paragraph{Gender.} The EEC operationalises gender through ten pairs of binary
noun phrases (\emph{she}/\emph{he}, \emph{my daughter}/\emph{my son}, \ldots)
and through first names, which implicitly assigns a gender to each name. This
excludes non-binary identities, whose erasure through binary treatments of
gender is itself a documented harm of language technologies~\cite{dev2021}, and
many pairs combine gender with kinship roles (\emph{mother}/\emph{father},
\emph{wife}/\emph{husband}). Following Larson~\cite{larson2017}, who recommends making
the operationalisation of gender explicit, we read our gender results as
sensitivity to these particular gendered expressions, not as evidence about
gender in general, and not as evidence about non-binary people.

\paragraph{Race.} Race is operationalised through first names that the EEC
labels African-American or European-American, drawn from earlier
work~\cite{caliskan2017,bertrand2004}. Names are imperfect proxies.
Gaddis~\cite{gaddis2017} shows that the race perceived from names used in audit studies
varies between names, and that names given more often by mothers with less
formal education are more likely to be perceived as Black; a contrast between
two name lists can therefore also be a contrast in perceived social class. The
categories are, moreover, specific to the United States and to
English~\cite{field2021}. Our race results are thus contrasts between two
particular lists of names. For brevity we refer to the two contrasts as
``gender'' and ``race'' throughout, in this restricted sense.

\paragraph{Consequences for interpretation.} Three follow. First, a detected
effect shows sensitivity to these tokens; attributing it to gender or race
requires the proxy assumptions above. Second, under the original paired test we
find reversals of direction (race in one classifier, gender in two LLM
configurations), which suggest that a summary such as ``biased against group
X'' can depend on the template type; the classifier reversal does not survive a
cluster-robust test (\S\ref{sec:eec}), so such summaries should at least be
stated conditionally on the templates and the test used. Third, we
pair names within gender and therefore do not estimate intersectional effects,
although the EEC name lists cross race and gender.

% ----------------------------------------------------------------------------
\section{Measuring Identity Sensitivity Across Label Schemes}\label{sec:measures}
The 18 models differ in label scheme, so we need a score that is comparable
across them. For a model with classes $c$ and softmax output $P(c\mid x)$, we
place each class on a valence axis with weight $w_c\in[0,1]$: negative and
positive receive 0 and 1, a neutral class 0.5, and five-point scales are evenly
spaced. These weights only follow the order of the labels; TV, the primary
measure for multi-class models, does not depend on them. The \emph{expected valence} is $e(x)=\sum_c w_c P(c\mid x)$. For a
sentence $x$ and its counterfactual $x'$, in which only the person term is
replaced by its matched counterpart, we consider three quantities:
\begin{itemize}
  \item $\Delta_e = e(x)-e(x')$, a score difference on the $[0,1]$ scale, the
  analogue of the score differences compared in the original EEC study;
  \item $\Delta_{\ell} = \operatorname{logit} e(x) - \operatorname{logit} e(x')$,
  computed stably from the logits. For binary models it equals the difference
  in logit margins and does not compress near 0 and 1; for multi-class models
  it is the logit of a probability-weighted average, which does compress when
  the model is confident in any one class, including the neutral class;
  \item $\mathrm{TV}(x,x') = \tfrac12 \sum_c |P(c\mid x) - P(c\mid x')|$, the
  total-variation distance between the two output distributions.
\end{itemize}
$\Delta_e$ and $\Delta_\ell$ are signed and can be used for directional tests.
Both, however, have a blind spot in multi-class models: if an identity swap
moves probability from the neutral class equally to the positive and negative
classes, $e$ is unchanged although the prediction has changed. TV responds to
any change in the output distribution but is unsigned. We therefore use
$|\Delta_\ell|$ as the primary sensitivity measure for binary models and TV for
multi-class models, and report the alternatives alongside (\S\ref{sec:measure}).
For binary models TV equals $|\Delta_e|$.

As a measure of predictive uncertainty we use $u(x) = 1-\max_c P(c\mid x)$.

% ----------------------------------------------------------------------------
\section{Experimental Setup}\label{sec:exp}
\paragraph{Models.} We audit 18 publicly available sentiment classifiers from
the Hugging Face Hub: 3 binary, 12 three-class and 3 five-class (including two
star-rating models), covering general-domain, social-media, review and
financial training data, and BERT, DistilBERT, RoBERTa and XLM-RoBERTa
architectures from base to large size. Full identifiers are listed in
Table~\ref{tab:models} in the appendix. All models are evaluated in English only.

\paragraph{Instruction-tuned LLMs.} We additionally audit two open
instruction-tuned language models, Qwen2.5-0.5B-Instruct and
Qwen2.5-1.5B-Instruct, used as zero-shot classifiers. Each sentence is placed in
the model's chat template with the instruction to answer with exactly one word,
either \emph{positive or negative} (binary prompt) or \emph{positive, neutral,
or negative} (three-class prompt); the class logits are the next-token logits of
the label words. Running the same weights under both prompts varies the label
scheme while holding the model fixed, which the classifier fleet cannot do. To
test sensitivity to wording, each scheme is run with the original prompt and
two alternatives that differ in structure, phrasing and answer-option order;
the alternatives were fixed, and their SHA-256 hashes recorded, before any of
them was run, and all are reported. The alternatives are run on the
minimal-pair probe only. For
CPU budget reasons, the minimal-pair probe for LLMs uses four emotion words (one
per emotion) instead of twenty and omits the control-token analysis; the EEC
replication and Study~3 are run in full.

\paragraph{Study 1: EEC replication.} We replicate the EEC template and pairing
protocol~\cite{kiritchenko2018}, adapting the output measures to the available
sentiment classifiers: the original systems were emotion-intensity regressors
scored on a $[0,1]$ scale, whereas we compare scores on the expected-valence and
logit scales (\S\ref{sec:measures}). We therefore replicate the audit design, not
the original systems or their results. For gender, each template--emotion-word
instantiation contributes eleven pairs: ten noun-phrase pairs (e.g.\ \emph{my
daughter}/\emph{my son}) and one pair of averages over female and male first
names. For race, each instantiation contributes one pair: the average over
African-American names against the average over European-American names. We
run a paired $t$-test per model on four subsets: the full EEC
(1,584 gender / 144 race pairs), its emotional templates (1,540 / 140), its
four neutral templates (44 / 4), and EEC-Neutral (176 / 16). EEC-Neutral
consists of the original four neutral templates plus twelve new ones of the same
kind (e.g.\ \emph{``$\langle$person$\rangle$ ordered a coffee.''}), filled with
the 60 EEC person terms using the EEC's own subject and object forms. Following
the original study, significance is Bonferroni-corrected over all assessments
(18 models $\times$ 2 axes: $\alpha = 0.05/36$). We report the mean absolute
difference $|\Delta_\ell|$ and the paired effect size $d_z$ (mean difference
divided by its standard deviation); unlike significance, $d_z$ does not depend
on the number of pairs.

\paragraph{Study 2: minimal pairs.} The EEC compares neutral and emotional
sentences built from \emph{different} templates. To compare emotional content
against a structurally matched filler condition, we build every sentence from
the same 16 base templates in three conditions:
\begin{quote}\small
\emph{bare:} I saw $\langle$person$\rangle$ in the market.\\
\emph{filler:} I saw $\langle$person$\rangle$ in the market, and then I went home.\\
\emph{emotion:} I saw $\langle$person$\rangle$ in the market, and then I felt
$\langle$emotion$\rangle$.
\end{quote}
A second clause form (\emph{``, which was on a Tuesday.''} / \emph{``, which
made me feel $\langle$emotion$\rangle$.''}) is used in the same way. Each
filler clause matches its emotion clause in position and syntactic form but not
word for word (\emph{went home} vs.\ \emph{felt $\langle$emotion$\rangle$}), so
\emph{filler vs.\ emotion} is the controlled contrast for emotional content
rather than for a single word; \emph{bare vs.\ emotion} corresponds to the
comparison implicit in the EEC and mixes emotional content with the added
clause. Emotion words are the EEC's 20
emotional-state words. Person terms are 5 gender noun-phrase pairs and 10
same-gender name pairs across race, chosen so that substitution never requires
other changes to the sentence, for 20,640 sentences per model. We compare mean
identity sensitivity between conditions as a ratio, with 95\% confidence
intervals from a bootstrap that resamples base templates and person pairs
independently (4,000 replicates), since both are random factors. Each pair is
counted once. A model is classed as showing an effect or a reversal only if the
two-sided bootstrap $p$-value for a ratio of 1 remains below 0.05 after Holm's
correction across the 18 models; the correction is applied separately for each
of the three measures, and all three are reported for every model in the
released results.

\paragraph{Study 3: uncertainty-aware template selection.} Because
probability-scale identity sensitivity is mechanically constrained by
predictive uncertainty (\S\ref{sec:unc}), the association between the two cannot by itself show that
uncertainty is useful for auditing. A consistent \emph{signed} shift against one
group, which is what the EEC test detects, is not mechanically implied by
uncertainty. We therefore turn the premise into a procedure and test a
prediction that can fail: template sets chosen to make a model uncertain should
detect directional bias more often than standard template sets of the same
size. The test uses $\Delta_\ell$, which for multi-class models also compresses
where the model is confident (\S\ref{sec:measures}), so for them the prediction
follows from the same mechanism; for binary models, whose $\Delta_\ell$ does
not compress, a gain would need another explanation. The procedure is
summarised below.

\begin{center}\fbox{\parbox{0.92\linewidth}{\small
\textbf{Uncertainty-aware template selection.} Given a model $f$, a pool of
candidate templates, a set of \emph{selection} person terms and a disjoint set
of \emph{test} person terms:
\begin{enumerate}
  \item fill every candidate template with each selection term and compute the
  mean predictive uncertainty $\bar u = 1-\max_c P(c\mid x)$ per template;
  \item keep the $n$ templates with the highest $\bar u$;
  \item run the standard counterfactual test on these templates with the test
  terms only, alongside the usual template sets;
  \item as a control, run the same test on within-group (placebo) swaps.
\end{enumerate}
Selection never uses an identity contrast or the test terms.}}\end{center}

Here, the pool has 320 frames, built by appending one of 20 mildly valenced or
factual clauses (e.g.\ \emph{``, and it was okay.''}, \emph{``, and it was a bit
dull.''}, \emph{``, and it was raining.''}) to each of the 16 base templates.
For each model, the \emph{targeted} set consists of the 16 frames with the
highest mean uncertainty, measured with 20 EEC first names that are never used
for testing. We compare it with three sets of 16 frames: a
\emph{random} draw from the same pool (to separate targeting from the clauses
themselves), the EEC-Neutral templates, and the base templates with an emotion
clause. Each set is tested with the same person pairs as Study~2 (80 gender and
160 race pairs per set), using the paired $t$-test on $\Delta_\ell$ with the
EEC's Bonferroni threshold. As a placebo, we swap names \emph{within} the same
race and gender group (256 pairs per set), where no directional effect is
expected; a set that turned noise into significance would flag placebo pairs.

\paragraph{Study 4: unseen template families.} To test whether the
uncertainty account carries over to sentence structures it was not developed
on, we wrote eight new template families whose structures do not occur in
Studies 1--3, for example a sentence-initial participle (\emph{Feeling furious,
I called my sister} / \emph{Being early, I called my sister}), an exclamative, a
question, and a family in which the person rather than the speaker experiences
the emotion (Table~\ref{tab:families} in the appendix). Each family has eight
base sentences, a matched filler, and, per emotion, the EEC word with the
highest and the lowest NRC intensity. Hypotheses, families, measures, decision
rules and a minimum number of informative cells were fixed and hashed before
any model scored the new sentences, and that minimum was checked before any
outcome was computed. Each model's uncertainty was measured with the 20
held-out names, and outcomes with the 15 Study~2 pairs. The unit is a model,
family and intensity level (288 cells). A cell's direction counts as clear if
the bootstrap interval of its log sensitivity ratio excludes zero; cells too
close to zero are not counted either way. H1 asks whether the uncertainty ratio
predicts the direction in clear cells (supported if the lower bound of a
two-way cluster-bootstrap interval over models and families exceeds 0.5); H2
asks the same for cells in which it disagrees with the model's Study~2
direction, and requires this on both the primary and a logit-scale outcome.

\paragraph{Label flips.} To relate sensitivity to decisions, we report the
share of counterfactual pairs whose predicted label (the most probable class)
changes under the identity swap.

\paragraph{Control tokens.} To test whether the emotion clause specifically
reduces the influence of the identity term, we compute deletion attributions on
the logit scale for a subset of sentences (one word per emotion) and compare the
identity term with the other tokens of the base template, which occur unchanged
in every condition. \emph{Specificity} is the factor by which the identity
term's attribution falls from filler to emotion, divided by the same factor for
the control tokens; a value of 1 means the emotion clause reduces the influence of
all context tokens alike.

\paragraph{Statistical inference.} We correct for multiplicity within each
family of tests, as follows. (i)~Directional tests in the EEC replication and in
Study~3, one paired $t$-test on $\Delta_\ell$ per model, identity axis and
template set, use the EEC's Bonferroni threshold $\alpha=0.05/36$ (18 models
$\times$ 2 axes); the placebo tests in Study~3 use the same threshold.
(ii)~Minimal-pair ratios are classified by Holm-corrected bootstrap $p$-values,
across the 18 classifiers separately for each output measure, and across the 12
LLM variants. (iii)~Model-level comparisons of which template set flags a bias
(Study~3) use exact McNemar tests with Holm correction over the six comparisons
(2 axes $\times$ 3 reference sets). (iv)~Wilcoxon signed-rank comparisons of
effect sizes and all correlations are reported descriptively, without
multiplicity correction. (v)~Because pairs in Studies~1 and~3 share templates
and person terms, we also report cluster-robust versions of the directional
tests, added after review: the mean $\Delta_\ell$ is tested with
cluster-robust standard errors (templates, and templates and person pairs, as
clusters; $t$ with $G-1$ degrees of freedom for the smaller number of
clusters $G$), at the same threshold, in line with calls for significance
tests that match the data's structure and for attention to statistical power
in NLP~\cite{dror2018,card2020}. The released code computes every reported
number from the saved results.

% ----------------------------------------------------------------------------
\section{Results}
Before the details, an overview of what holds and what does not.
\emph{Robust:} the same template change moves identity sensitivity in opposite
directions in different models, closely following uncertainty
(\S\ref{sec:mp}--\ref{sec:unc}), also on eight unseen template families in a
test fixed in advance, on the probability scale on which audits register effects
(\S\ref{sec:families}), and audit verdicts change with the type and
number of templates, with the output measure and with the significance test
(\S\ref{sec:eec}, \S\ref{sec:measure}). \emph{Conditional:} uncertainty-aware
templates flag disparities in more models under the original paired test, but
not under a cluster-robust one, and even there significantly only against emotional
templates, and a cross-model control cannot attribute the gain to selecting
templates for the audited model (\S\ref{sec:targeted}); whether uncertainty
adds information beyond a model's usual direction depends on the outcome
measure (\S\ref{sec:families}); and the LLM evidence rests on two models from
one family (\S\ref{sec:llm}). \emph{Negative:} in natural review sentences,
emotional content barely changes uncertainty, so the template contrast has
little counterpart there and could be neither confirmed nor refuted
(\S\ref{sec:natural}).

\subsection{RQ1: The EEC pattern in the EEC itself}\label{sec:eec}
Table~\ref{tab:eec} summarises the EEC replication (per-model values in
Table~\ref{tab:eecfull}). The original pattern, a larger mean identity
difference on neutral than on emotional templates, holds for 7 of 18 models on
gender and 5 of 18 on race; the median difference is in fact smaller on
EEC-Neutral (gender 0.068 vs.\ 0.122; race 0.036 vs.\ 0.122 on the logit
scale). Effect sizes show the opposite tendency: the median $|d_z|$ is larger on
EEC-Neutral (gender 0.35 vs.\ 0.16; race 0.59 vs.\ 0.40), and among models with
a significant full-EEC bias, $|d_z|$ is larger on EEC-Neutral in 8 of 13
(gender) and 7 of 11 (race). On emotion-free text, identity-related shifts are
thus typically smaller but more consistent in direction. Across models, the
EEC neutral-to-emotional ratio of $|\Delta_\ell|$ correlates with the
minimal-pair TV ratio of \S\ref{sec:mp} (Spearman $\rho=0.64$, $p=0.004$, gender), so the two
instruments agree on which models show the pattern.

\begin{table}[t]\centering\footnotesize
\caption{EEC replication: number of models (of 18) with a significant bias on
$\Delta_\ell$, Bonferroni $\alpha=0.05/36$. \emph{t}: the paired $t$-test of the
original protocol, which treats every pair as independent. \emph{Robust}: a
cluster-robust test with templates and person pairs as sources of variation
(for race, each comparison is already averaged over names, so only templates
apply). Results of the $t$-test on the $[0,1]$ scale ($\Delta_e$) are similar:
gender 12/6/9 and race 9/0/6 for the first three rows. Lower block (original
test only): emotional pairs subsampled to the EEC-Neutral size (2,000 draws per
model); ``usually'' means in more than half of the draws. The seven misses not
usually found by an equal-sized emotional set are found in at most 20\% of
draws.}
\label{tab:eec}
\setlength{\tabcolsep}{4pt}
\begin{tabular}{lcccc}
\toprule
& \multicolumn{2}{c}{Gender} & \multicolumn{2}{c}{Race} \\
\cmidrule(lr){2-3}\cmidrule(lr){4-5}
& \emph{t} & Robust & \emph{t} & Robust \\
\midrule
Significant on full EEC & 13 & 2 & 11 & 6 \\
Significant on original 4 neutral templates (44 / 4 pairs) & 6 & 0 & 0 & 0 \\
Significant on EEC-Neutral (176 / 16 pairs) & 10 & 1 & 6 & 6 \\
\midrule
Detected by EEC-Neutral, missed by original neutral templates & 4 & 1 & 6 & 6 \\
Significant on full EEC, not on EEC-Neutral & 4 & 2 & 5 & 3 \\
Significant on EEC-Neutral, not on full EEC & 1 & 1 & 0 & 3 \\
Significant on both, opposite directions & 0 & 0 & 1 & 0 \\
\midrule
Mean $|\Delta_\ell|$ larger on EEC-Neutral than emotional & \multicolumn{2}{c}{7} & \multicolumn{2}{c}{5} \\
\midrule
Expected flags, emotional set of EEC-Neutral size & 6.9 & -- & 3.0 & -- \\
Missed by EEC-Neutral, usually found by equal-sized emotional set & 2 & -- & 0 & -- \\
Flagged by EEC-Neutral, usually missed by equal-sized emotional set & 5 & -- & 4 & -- \\
\bottomrule
\end{tabular}
\end{table}

Four points stand out. First, how the test treats identity terms changes the
gender verdicts more than any other choice. The paired $t$-test of the original
protocol treats every pair as independent, although pairs share templates and
person terms. With templates and person pairs as sources of variation
(Table~\ref{tab:eec}, \emph{Robust}), the number of models with a significant
gender bias on the full EEC falls from 13 to 2 (9 if only templates are
treated as varying): the EEC's gender contrast rests on ten noun-phrase pairs
and one average over names, and its directional gender verdicts do not
generalize beyond those pairs. Race comparisons, already averaged over names,
are less affected (11 to 6). Second, the original neutral subset is too small
to be informative: its four race pairs never reach significance, and
EEC-Neutral recovers 4 gender and 6 race biases that it misses (1 and 6 under
the robust test), including the largest race bias in the fleet ($d_z = -2.89$
on EEC-Neutral). Third, both the type and the number of templates shape the
verdict. Under the original test, 4 gender and 5 race biases are
significant on the full EEC but not on EEC-Neutral, which is much smaller (176
gender and 16 race pairs, against 1,540 and 140 emotional pairs). To separate
type from size, we drew for each model 2,000 random subsets of emotional pairs
of the EEC-Neutral size (Table~\ref{tab:eec}, lower block). An emotional set of
that size would usually detect only two of the nine (both gender, in 80\% and
100\% of draws); the other seven would usually be missed at that size as well,
so most of these misses reflect the small number of neutral pairs. At equal
size, type matters more often in the other direction: EEC-Neutral flags 5
gender and 4 race biases that an equal-sized emotional set would usually miss.
Under the robust test, race verdicts still depend on the template type: 3 race
biases are significant only on the full EEC and 3 only on EEC-Neutral. Either
way, an audit that replaced emotional templates with these neutral ones would
pass the missed models. Fourth, under the original test, the race bias of one
model (twitter-roberta) is significant on both template types with opposite
signs: African-American names receive higher scores on emotional sentences
($d_z=+0.52$) and lower scores on neutral sentences ($d_z=-1.15$), and because
the full EEC is 97\% emotional pairs, the standard test reports only the
first. Under the robust test, the emotional-template side is no longer
significant, so this reversal is not robust.

Finally, the binary models illustrate the gap between sensitivity and
directional bias: their identity differences on EEC-Neutral are larger than on
emotional templates (by factors of 1.5--7.3 for gender), but for two of the
three the direction is not consistent across pairs and no gender bias is
significant on neutral templates; only distilbert-sst-2 shows a significant
gender bias there ($d_z=-0.51$).

\subsection{RQ2: Replacing the emotion clause has model-dependent effects}\label{sec:mp}
Table~\ref{tab:mp} reports, for each model, the ratio of identity sensitivity in
the filler condition to the emotion condition. A ratio above 1 means the
identity term matters more in the filler than in the emotion condition, the direction suggested by
the EEC observation. Under the primary measure and after Holm correction, the
ratio is significantly above 1 for 5 models, not distinguishable from 1 for 8,
and significantly below 1 for 5. All three binary models show the effect
(ratios 3.4--10.5 on the logit scale), as do one star-rating model and one
three-class model. All three
financial-domain models, and the large three-class RoBERTa model, show strong
effects in the opposite direction: on emotion-free sentences, identity
substitution changes their output very little.

The less controlled \emph{bare vs.\ emotion} contrast points toward the effect
more often (10 above 1, 4 below; Table~\ref{tab:mp}). The difference arises
because adding a neutral clause alone reduces identity sensitivity in several
models; a bare-vs-emotion comparison, like the EEC's, therefore attributes to
emotional content part of an effect that is due to sentence form.

\begin{table}[t]\centering\scriptsize
\caption{Minimal-pair results per model, sorted by the primary ratio. $k$:
number of classes. $P_{\text{neu}}$: mean probability of the neutral or middle
class on filler sentences. $\bar u$: mean predictive uncertainty
$1-\max_c P(c\mid x)$ on filler and emotion sentences. Primary ratio: mean
identity sensitivity, filler over emotion, on $|\Delta_\ell|$ for binary models
and TV for multi-class models, with 95\% two-way bootstrap CI; $\uparrow$
significantly above 1, $\downarrow$ significantly below 1, $\circ$ not
significant, all after Holm correction across the 18 models. Valence: the same ratio on
$|\Delta_\ell|$ for every model, with markers as for the primary ratio. Bare/emo: TV ratio, bare over emotion.
Spec.: control-token specificity (1 = no identity-specific effect).}
\label{tab:mp}
\setlength{\tabcolsep}{3pt}
\begin{tabular}{lcccc l c ccc}
\toprule
Model & $k$ & $P_{\text{neu}}$ & $\bar u_{\text{fill}}$ & $\bar u_{\text{emo}}$ &
Primary ratio [95\% CI] & & Valence & Bare/emo & Spec. \\
\midrule
\inputrows{figures_tables/table_minimal_pair.tex}
\bottomrule
\end{tabular}
\end{table}

\paragraph{The effect is not specific to identity.} Median control-token
specificity is 0.85: the emotion clause reduces the influence of the identity term
no more than that of the other template tokens (specificity significantly above
1 in 3 models, below 1 in 8, and not distinguishable from 1 in 7). Consistently,
the identity term's share of total attribution is higher in filler than in
emotion sentences in 16 of 18 models (mean 0.112 vs.\ 0.053, with the same
uniform baseline of 0.097 in both conditions because the clauses have equal
length), but this reflects the emotion clause taking attribution from all tokens,
not from the identity term in particular. We therefore do not interpret the
results as the emotion clause ``crowding out'' identity information.

\subsection{RQ3: Identity sensitivity follows predictive uncertainty}\label{sec:unc}
The models that reverse share a profile visible in Table~\ref{tab:mp}: they
place emotion-free sentences in the neutral class with near certainty
($P_{\text{neu}} \geq 0.87$ for the four strongest reversals), while emotion
words move them only partly off neutral, leaving emotional sentences in a
region of higher uncertainty. Binary models have no neutral class and show the
opposite profile: they are confident on emotional sentences and less certain on
emotion-free ones.

\begin{figure}[t]\centering
\includegraphics[width=0.82\linewidth]{figures_tables/fig_uncertainty.pdf}
\caption{Identity-sensitivity ratio against uncertainty ratio (filler over
emotion sentences), one point per model, on the primary measure
($|\Delta_\ell|$ for binary, TV for multi-class models; bars: 95\% CI).
Points right of the vertical line are less certain in the filler condition;
points above the horizontal line are more sensitive to the identity term in
the filler condition. Hollow markers are not significant after Holm correction.}
\label{fig:unc}
\end{figure}

Across the 18 models, the log ratio of identity sensitivity (filler over
emotion, measured by TV, which is defined in the same way for every label
scheme, so no model-specific choice of measure enters) is strongly correlated with the log ratio of predictive uncertainty
in the same conditions (Spearman $\rho = 0.91$, $p<10^{-6}$; $\rho=0.90$ on the
primary measure, Fig.~\ref{fig:unc}); the same holds for bare over emotion
($\rho=0.89$). The association is not driven by individual models: leaving out
any one model, it ranges from 0.89 to 0.95 (0.88 to 0.94 on the primary
measure), and a bootstrap over models gives a 95\% interval of 0.69--0.98
(0.68--0.97). Within every model, sentence-level uncertainty
and identity sensitivity are positively correlated (median $\rho=0.60$, range
0.27--0.94 across models). Neutral-class probability alone is a weaker
predictor: among the 15 multi-class models it correlates with the sensitivity
ratio at $\rho=-0.70$, compared with $\rho=0.85$ for the uncertainty ratio. One
model with little neutral mass (amazon-review, 0.09) still points towards a
reversal (0.72), although not significantly after Holm correction, because it
assigns emotion-free sentences fairly confidently to a positive rating.

We stress that this association is partly expected by construction: near a
probability of 1 the softmax is flat, so a small change in a model's logits
moves its probabilities, and hence TV, very little. The association also holds
on $|\Delta_\ell|$ (Spearman $\rho=0.85$ across the 18 models, using
$|\Delta_\ell|$ for every model), but this removes the coupling only for the
three binary models: for multi-class models, $\Delta_\ell$ is the logit of
expected valence, which also compresses when the model is confident in any one
class (\S\ref{sec:measures}). Nor does it
depend on how uncertainty is defined: with normalised entropy or with one minus
the margin between the two most probable classes, the correlation is 0.91 under
both definitions on TV (0.84 and 0.87 on $|\Delta_\ell|$), and the median
within-model correlation is 0.61 and 0.60. Under entropy, three models whose
uncertainty ratios are close to 1 change direction, one of them a significant
reversal (BERT\_sentiment).

\paragraph{How much of the direction is saturation?} In an exploratory check,
we divided each model's logits by a temperature $T$, which leaves the predicted
class unchanged but removes saturation as $T$ grows. For 14 of the 18 models,
including 8 of the 10 with a significant effect, the probability-scale direction
at $T=1$ already matches the direction of the change in the model's logits and
does not change with $T$. For the other 4, among them the reversal of
BERT\_sentiment and the effect of bert-multi-5star, it flips as saturation is
removed, and at $T=16$ all 18 match the logit-space direction. The
association with uncertainty weakens but persists ($\rho=0.73$ at $T=8$ and
$T=16$). Saturation thus decides the direction for a minority of models; for
most, the probability-scale direction reflects how the model's logits respond.

\paragraph{The association predicts held-out person terms.} Uncertainty and
sensitivity above are measured on the same sentences. In an exploratory check
added after the main analyses, we split the 15 person pairs at random into two
halves, measured each model's uncertainty ratio on one half and its sensitivity
ratio on the other, and repeated this 200 times. The correlation across the 18
models remains high (median $\rho=0.88$, 5th percentile 0.79; on
$|\Delta_\ell|$, 0.85), and the uncertainty ratio from one half gives the right
direction of the sensitivity ratio on the other for all 10 models with a
significant effect in every split (9 of 10 on $|\Delta_\ell|$). The prediction
holds across person terms within the same templates and models; it does not
show prediction across templates, tasks or new models.

The finding is therefore not that uncertainty causes bias or harm, but that identity
effects are \emph{observed} mainly where a model is uncertain, and that which
templates make a model uncertain differs systematically between models.

\paragraph{The effect reaches decisions.} The same pattern appears in predicted
labels. Identity swaps change the predicted label more often on emotional than
on filler sentences in 13 models and less often in 4 (the three binary models
and bert-multi-5star). Across models, the ratio of flip rates follows the
sensitivity ratio (Spearman $\rho=0.71$, $p<0.001$). Flip rates are not
negligible where the model is uncertain: up to 10.8\% of filler pairs for the
binary Roberta\_Sentiment and 16.0\% for bert-multi-5star, and 8.1--8.2\% of
emotion pairs for finbert-tone and amazon-review, compared with near zero in the
condition where each model is confident.

\subsection{A prospective test on unseen template families}\label{sec:families}
Study~4 (\S\ref{sec:exp}) tests the account on eight template families it was
not developed on. Of the 288 cells (\emph{model} $\times$ \emph{family}
$\times$ \emph{intensity}), the direction is clear in 168 cells. The other 120
cells are too close to zero to assign a direction and are not counted either
way.

\paragraph{The account generalizes to the new families.} The uncertainty ratio
gives the right direction in 93\% of clear cells (95\% CI from a bootstrap over
models and families, 0.82--0.99), and it correlates with the sensitivity ratio
across all 288 cells at $\rho=0.88$ (0.73--0.94). Leaving out any one family,
the accuracy stays between 0.91 and 0.95. Rules that assign each model a fixed
direction do worse on the same cells: the model's own Study~2 direction is right
in 0.74, the label scheme in 0.67 and neutral-class probability in 0.68.

\paragraph{It predicts departures from a model's usual direction, on the
primary measure only.} In 98 cells, counted before any outcome was computed, the
uncertainty ratio pointed the opposite way to the model's Study~2 direction. In
the 45 of these with a clear outcome it was right in 84\% (0.56--1.00), and
within models the two ratios correlate at 0.77 on average. On the logit-scale
outcome, which we included as a guard against mathematical coupling, this
advantage is not established (0.73, 0.37--0.95), and the label scheme predicts
the direction better than uncertainty does (0.90 against 0.79 over clear
cells). By the rules fixed in advance, uncertainty therefore adds information
beyond a model's usual direction on the primary measure but not on the logit
outcome. An entropy-based measure of uncertainty gives the same results.

\paragraph{Exploratory follow-up.} Checks added after the frozen analysis
(Appendix~\ref{app:families}) show that excluding cells without a clear
direction does not create the effect, that the two outcomes differ only in
multi-class models, and that for those models a measure that does not saturate
is predicted no better by uncertainty than by each model's usual direction.

\subsection{The output measure can reverse the verdict}\label{sec:measure}
For multi-class models, the valence-based measure and TV can disagree in
direction (Table~\ref{tab:mp}, \emph{Valence} vs.\ \emph{Primary ratio}). On
$|\Delta_\ell|$, 13 of the 15 multi-class models appear significantly reversed
after Holm correction; under TV only 5 do, and for twitter-xlm-roberta the
conclusion changes from reversed (0.53) to a significant effect (1.83). The discrepancy arises because these
models often hold most probability on the neutral class, and identity swaps move
probability out of it toward both other classes, which leaves the expected
valence nearly unchanged. Conversely, for binary models ratios on the $[0,1]$
scale are unstable when emotional predictions are saturated (for the large
binary RoBERTa, the TV ratio is 16.0 with CI [0.72, 240]), whereas the
logit-scale ratio is well defined (3.65 [1.53, 7.31]).

\subsection{Study 3: Uncertainty-aware template selection}\label{sec:targeted}
Table~\ref{tab:targeted} compares the four template sets. Targeted sets reach
the highest uncertainty in every model and the highest median uncertainty
overall (0.42--0.47 against 0.11--0.20). They flag a significant directional
gender disparity in 10 of 18 models, against 6 for EEC-Neutral, 5 for random frames and 2 for
emotional templates, and a race bias in 8, tied with emotional templates. No set
flags any placebo comparison, so the additional flags are not an artefact
of amplified noise. Compared with random frames from the same pool, targeted
sets flag 6 gender and 3 race biases that random sets miss, against 1 and 1 in
the other direction, which suggests that targeting rather than the clause pool
drives the difference.

These counts are descriptive. Comparing, model by model, which template set
flags a bias (exact McNemar test on the discordant models, Holm correction over
the six comparisons), the only difference that remains significant is targeted
versus emotional templates for gender (8 vs.\ 0 discordant models,
$p_{\text{Holm}}=0.047$). Against random frames and EEC-Neutral, the discordant
models favour targeted sets for gender (6 vs.\ 1 and 5 vs.\ 1), but with 18
models these differences are not significant ($p=0.125$ and $0.219$ before
correction); for race, no comparison approaches significance. The evidence for a
detection advantage is therefore consistent in direction but modest in strength.

\paragraph{The detections do not survive a cluster-robust test.} The counts
above use the paired $t$-test, which treats the 80 gender pairs of a set (16
frames $\times$ 5 noun-phrase pairs) as independent. Treating only frames as a
source of variation is not enough: it lowers the targeted counts to 6 gender
and 5 race detections, but it also flags placebo swaps in 4 models for the
targeted set, 4 for EEC-Neutral and 3 for emotional templates, because
individual names differ systematically. With frames and person pairs both as
sources of variation, no template set flags any gender, race or placebo
comparison in any model (Table~\ref{tab:targeted}). With five gender and ten
race pairs per set, Study~3 therefore cannot show directional disparities that
generalize beyond the specific identity terms, and the detection advantage
above should be read as descriptive of these terms only. Where a targeted set and the full EEC both find a significant bias,
they agree in direction in 12 of 13 cases (gender 7/7, race 5/6; the exception
is the race bias of distilbert-sst-2, consistent with the context-dependence of
direction reported in \S\ref{sec:eec}).

Effect sizes show a similarly uneven picture: across models, targeted $|d_z|$
is significantly larger in three of the six comparisons (without multiplicity
correction) and not in the others.\footnote{Wilcoxon signed-rank tests across models, not corrected for
multiplicity: targeted $|d_z|$ is larger than that of random frames for race
(Wilcoxon signed-rank, $p=0.027$; gender $p=0.054$), of EEC-Neutral for race
($p=0.006$) and of emotional templates for gender ($p=0.048$).} Targeted sets also miss 1 gender and 3 race biases
that neutral or emotional templates flag, so they complement rather than
replace standard templates. For the three models that are near-certain almost
everywhere (FinancialBERT, finbert-tone, roberta-large-3cls), targeting raises
uncertainty from below 0.01 to 0.19--0.30, yet no set finds a significant bias;
targeting increases the chance of observing identity effects but cannot create
an effect that is absent.

\paragraph{Targeting or generally uncertain frames.} Because the targeted
frames come from a pool of valenced clauses, the gain could come from frames
that make every model uncertain rather than from selecting frames for the
audited model. As a control, fixed with a recorded hash before scoring, we
tested each model on the 16 frames with the best mean uncertainty rank across
the other 17 models. These consensus sets flag a gender disparity in 7 models
and a race disparity in 8 (targeted sets: 10 and 8; no placebo detections in
either), and applying another model's targeted set flags on average 5.9 gender
and 8.4 race disparities. For gender the discordant models favour targeted sets
(5 vs.\ 2), but not significantly (exact McNemar, $p=0.45$); for race there is
no difference (4 vs.\ 4). The control therefore does not establish that the
gender gain requires selecting templates for the audited model. It does show
that uncertainty is strongly model-specific: targeted sets of different models
barely overlap (mean Jaccard index 0.03), consensus frames leave the audited
model about half as uncertain as its own targeted frames (median 0.21 against
0.47), and the targeted sets draw evenly on mildly negative, mildly positive and
factual clauses (102, 89 and 97 of 288 frames), so no single kind of clause
drives the selection.

The advantage does not clearly replicate in the four LLM variants (original
wordings). There,
targeted sets flag a gender bias in 3 variants (EEC-Neutral 2, emotional 2,
random 1) but a race bias in only 1 (EEC-Neutral 2, random 2, emotional 0),
again with no placebo detections. Pooled over the 18 classifiers and four LLM
configurations, targeted sets
detect the most gender biases (13, against 8 for EEC-Neutral, 6 for random and 4
for emotional templates) but are essentially tied for race (9, against 8 for
random, 8 for emotional and 7 for EEC-Neutral). We therefore regard
uncertainty-aware template selection as a diagnostic that complements standard
template sets, most clearly for gender, rather than as a validated replacement.

\begin{table}[t]\centering\footnotesize
\caption{Study 3: number of models (of 18) with a significant directional bias
per template set (16 frames each; Bonferroni $\alpha=0.05/36$), median $|d_z|$,
and median predictive uncertainty $\bar u$. \emph{t}: paired $t$-test on
$\Delta_\ell$; \emph{Robust}: cluster-robust test with frames and person pairs
as sources of variation (gender / race / placebo). Placebo: name swaps within
the same race and gender group.}
\label{tab:targeted}
\setlength{\tabcolsep}{3pt}
\begin{tabular}{l ccc c ccc c}
\toprule
& \multicolumn{3}{c}{Models flagged (\emph{t})} & Robust & \multicolumn{3}{c}{Median $|d_z|$} & \\
\cmidrule(lr){2-4}\cmidrule(lr){6-8}
Template set & Gender & Race & Placebo & G / R / P & Gender & Race & Placebo & $\bar u$ \\
\midrule
Targeted (max.\ uncertainty) & \textbf{10} & \textbf{8} & 0 & 0 / 0 / 0 & \textbf{0.40} & \textbf{0.25} & 0.05 & 0.42--0.47 \\
Random (same pool) & 5 & 6 & 0 & 0 / 0 / 0 & 0.22 & 0.20 & 0.04 & 0.20 \\
EEC-Neutral & 6 & 5 & 0 & 0 / 0 / 0 & 0.26 & 0.17 & 0.03 & 0.11--0.14 \\
Emotional & 2 & \textbf{8} & 0 & 0 / 0 / 0 & 0.18 & 0.24 & 0.05 & 0.12--0.13 \\
\bottomrule
\end{tabular}
\end{table}

\subsection{A bias of known size}\label{sec:injection}
Whether an audit can detect a bias, rather than only how sensitive a model is,
cannot be checked without a bias of known size. In an exploratory simulation,
we injected a directional bias into every female sentence of the EEC by
shifting the model's logits towards positive by $\delta$ (for binary models,
$\delta$ is added to the logit margin), and tested the gender axis with the
cluster-robust test. On the probability scale used by the original EEC study,
an injected bias of $\delta=0.5$ was detected on emotional templates in 7 of
the 18 models, against 13 on the logit scale; at $\delta=1$, in 11 against 15.
Across models and template sets, the size of the injected effect on the
probability scale follows the set's mean uncertainty ($\rho=0.84$ at
$\delta=0.5$): the same bias looks smaller where the model is confident. This
shows the premise of the paper on an effect of known size. Comparisons between
template sets are confounded by their size and are not interpreted.

\subsection{Instruction-tuned LLMs and a controlled label-scheme test}\label{sec:llm}
Table~\ref{tab:llm} in the appendix reports the minimal-pair results for the LLMs under three
prompt wordings per label scheme. The two alternative wordings (p2, p3) differ
from the original (p1) in structure, phrasing and the order of the answer
options; they were fixed before any of them was run, and all are reported.

\paragraph{The direction is robust to wording; the size is not.} In 11 of the
12 variants, identity sensitivity is significantly higher on filler than on
emotion sentences (Holm correction across the 12 variants). The magnitude,
however, depends strongly on the wording: among these 11, ratios range from 1.4
to 9.3, and the
original wording gives the largest ratio in three of the four model--scheme
combinations (binary prompts: 3.0--4.0 with p1, 1.4--2.7 with p2 and p3). We
therefore interpret the direction, not the size, of the LLM effect. Label flips
follow the sensitivity ratio less consistently than in the classifiers: in 3 of
the 12 variants, flips are more frequent on emotion than on filler sentences.

\paragraph{Uncertainty tracks the direction within one model.} In 10 of the
11 variants with a significant effect, filler sentences are also the less
certain condition; the exception (0.5B, three-class, p3) has nearly equal
uncertainty in both conditions (ratio 0.98). Across the 12 variants, the
sensitivity ratio correlates with the uncertainty ratio at Spearman
$\rho=0.66$ ($p=0.02$); the variants share weights and are not independent, so
we report this separately from the classifier analysis. With the original
prompts, pooling the 18 classifiers and the four LLM variants gives
$\rho=0.90$ ($p<10^{-7}$), and in all 14 systems with a significant effect or
reversal the direction matches that of the uncertainty ratio.

\paragraph{Changing the prompt wording moves a model between regimes.}
Under the original three-class prompt, Qwen2.5-1.5B assigns 89\% of probability
on filler sentences to \emph{neutral}, the behaviour of the neutral-locked
classifiers in Table~\ref{tab:mp}. Its expected-valence measure then shows an
apparent reversal (0.25), while TV shows a clear effect (9.27): the measurement
discrepancy of \S\ref{sec:measure}, with the weights held fixed. The direction
does not reverse, because emotional sentences remain more certain (uncertainty
0.008 vs.\ 0.064). Under the p3 wording, the same model becomes fully
neutral-locked (99.7\% neutral on filler sentences, uncertainty 0.003 against
0.034 on emotion sentences): it is the only variant whose uncertainty ratio
falls below 1, and the only one without a significant effect (0.81 [0.10,
1.88]; on the valence measure, 0.06). The p3 wording differs from the original
in phrasing, structure and the order of the answer options; changing only the
answer options, from the binary to the three-class prompt with the same wording,
leaves the effect in place for both models (Table~\ref{tab:llm}). With the
weights and sentences fixed, the effect thus loses significance in the one
variant whose wording leaves emotion-free text more certain than emotional
text. The account predicts a ratio below 1 there; the estimate is below 1 but
its interval is wide, so this is a single case consistent with the account,
not a confirmed reversal.



In the EEC (run with the original wordings), the raw identity difference is
larger on EEC-Neutral than on emotional templates for three of the four variants (1.5--2.9 times; the
exception is Qwen2.5-1.5B with the three-class prompt). Under the binary prompt,
both model sizes show a significant gender bias on emotional templates
($d_z=-0.21$ and $-0.46$) and a significant gender bias in the opposite
direction on EEC-Neutral ($+0.31$ and $+0.51$) under the original paired
test, a second instance of the direction reversal reported in
\S\ref{sec:eec}, with the same caveat about the test.

\subsection{A check on natural sentences}\label{sec:natural}
All results so far use synthetic templates. To test whether they carry over to
natural text, we took sentences from the unlabelled portion of the IMDB review
corpus~\cite{maas2011}, which is not part of the labelled IMDB data used to
fine-tune sentiment classifiers, although some checkpoints may have been trained
on other IMDB reviews,
and kept 1,000 sentences that refer to people only through gendered words of a
single gender (pronouns and words such as \emph{son}, \emph{wife} or
\emph{king}) and contain no names. Each sentence was given a full gender swap,
yielding a male and a female version. Half of the sentences contain no word from
the NRC Affect Intensity Lexicon~\cite{mohammad2018} (low affect) and half
contain a word with intensity of at least 0.5 (high affect), a natural-text
counterpart of the filler and emotion conditions. The sampling rules, swap rules
and tests were fixed, with a recorded hash, before any model was scored; after
reading sample swaps, and still before scoring, we tightened the name filter
and the word list. In a manual check of 30 sentences containing a female
\emph{her}, two or three conversions were wrong or doubtful (e.g.\ \emph{make
her crazy} $\to$ \emph{make his crazy}), which affects about 1\% of all
sentences. We scored the 18 classifiers and the four LLM configurations with the
original prompts (Table~\ref{tab:natural} in the appendix).

\paragraph{The within-model association replicates.} In all 22 systems,
sentence-level uncertainty and identity sensitivity (TV) are positively
correlated, significantly after Holm correction (median $\rho=0.74$, range
0.20--0.97), compared with a median of 0.60 on templates; with entropy or the
top-two margin as the uncertainty measure, the medians are 0.74 and 0.73, again
positive in all 22 systems. In real sentences, as
in templates, a gender swap moves the output mainly where the model is
uncertain, with the same mechanical caveat as in \S\ref{sec:unc}.

\paragraph{Affect barely changes uncertainty here, and the contrast is weak.}
Unlike the template contrast, the two strata differ little in uncertainty: the
uncertainty ratio (low over high affect) lies between 0.8 and 1.2 for 20 of the
22 systems. Only one system has a sensitivity ratio that is significant after
Holm correction: FinancialBERT, which is much less uncertain on low-affect
sentences (uncertainty ratio 0.36) and correspondingly less sensitive there
(0.41). Across systems, the two log ratios correlate at $\rho=0.34$ ($p=0.12$;
0.40, $p=0.10$, for the classifiers alone). This exceeds the threshold of 0.3
we fixed in advance, but it is not significant, so we count the
between-condition prediction as neither supported nor refuted in natural text:
with so little difference in uncertainty between the strata, the test has
little leverage. The between-condition result of Study~2 therefore rests on
templates built to differ in uncertainty.

\paragraph{Verdicts again depend on the sentences audited.} A directional
gender bias ($\alpha=0.05/44$, Bonferroni over 22 systems and two strata, by
analogy with the EEC threshold of 0.05/36 for 18 models and two axes) is
detected in 13 systems on low-affect
sentences and in 15 on high-affect sentences, and in 17 on at least one. For 6
of the 22 systems, one stratum detects a bias and the other does not; unlike
in the EEC, no direction reverses. The direction differs between systems: of
the detections, some favour the female and some the male versions, and all
four LLM configurations score male versions more positively in both strata.

% ----------------------------------------------------------------------------
\section{Discussion}
\paragraph{What the EEC observation reflects.} The larger neutral-template
differences reported for the original EEC systems also appear here in the
binary classifiers and a minority of multi-class models, but not in models that
confidently label emotion-free text as neutral. Our results are consistent with
an account in which identity effects are observable where the model's
prediction is not saturated: for models that are uncertain on plain text, that
is the neutral templates; for models with a dominant neutral class, it is the
emotional ones. We do not claim this explains the behaviour of the original
SemEval-2018 systems, which were emotion-intensity regressors rather than
polarity classifiers and are not publicly available.

\paragraph{What goes beyond saturation, and what does not.} Saturation alone
predicts only that identity effects are hard to see where a model is
confident. Several results go further. The same change of template moves
identity sensitivity in opposite directions in different models
(\S\ref{sec:mp}), as Selvam et al.~\cite{selvam2023} also observed for bias
measures, and the frames that make one model uncertain rarely make another
uncertain (\S\ref{sec:targeted}): the profiles that determine where effects are
visible are a property of each model. Each model's direction is also stable:
its Study~2 direction predicts most of its directions on unseen template
families, on every measure we used (\S\ref{sec:families}). And the
\emph{direction} of a detected bias, which uncertainty does not determine,
reverses between template types under the original paired test in one
classifier and two LLM configurations, although the classifier case does not
survive a cluster-robust test (\S\ref{sec:eec}, \S\ref{sec:llm}). What the
uncertainty measurements add is narrower. Measured on held-out person terms,
they predict the direction of probability-scale sensitivity on other terms
(\S\ref{sec:unc}) and on unseen template families, including where a model
departs from its usual direction (\S\ref{sec:families}). On log-scale measures
this incremental information is not established, and for multi-class models a
measure that does not saturate is predicted no better by uncertainty than by
the model's usual direction (\S\ref{sec:families}). The predictions therefore
concern where effects become visible on the probability scale on which audits
register them, which is the paper's premise, rather than a property of the
models' underlying response beyond saturation. The one LLM prompt that removed
the effect (\S\ref{sec:llm}) is consistent with this reading.

\paragraph{From observability to verdicts.} The predictive results concern
unsigned sensitivity; an audit verdict concerns a signed disparity. The
evidence connecting the two is weaker and should be read as such. Three
classifiers that are near-certain on emotion-free text show almost no identity
effect there and larger effects on emotional templates (Table~\ref{tab:eecfull}, ratio columns);
at equal size, neutral and emotional templates flag different disparities, but
only two of the nine disparities missed by EEC-Neutral are attributable to
template type (\S\ref{sec:eec}); and templates selected for high uncertainty
flag gender disparities in more models under the original paired test,
significantly so only against emotional templates, but no template set flags
any disparity under a cluster-robust test (\S\ref{sec:targeted}). With the few
identity pairs of these benchmarks, directional verdicts mainly describe the
specific terms tested. Low uncertainty on every template
thus marks an audit as uninformative, but high uncertainty does not guarantee
that a disparity, if present, will be flagged. Our results therefore show when
identity sensitivity is observable to a template audit; they do not show that
uncertainty-guided auditing is a better method for detecting discrimination.

\paragraph{What the results imply jointly.} An audit verdict is a joint
property of the model, the template set (its type and its size) and the output
measure. Which templates make a model uncertain differs sharply between models
(the 16 most uncertain frames of different models barely overlap,
\S\ref{sec:targeted}), so the same template set can be informative for one model and nearly
uninformative for another, and a standard audit report does not show which. On
natural review sentences, the within-model association held in every system,
but emotional content did not create the large differences in uncertainty that
emotion-bearing templates did (\S\ref{sec:natural}), so the conditions under
which a template audit observes a model need not match those of the text it
will process.

\paragraph{What remains open.} We cannot rule out that uncertainty partly marks
something else, such as how unfamiliar an input is to a model; the prompt
manipulation speaks against a purely textual explanation, but we have not
manipulated uncertainty independently of the model. Whether the pattern holds
for other tasks, languages and identity categories, and for uncertainty
estimated from ensembles, is untested.


% ----------------------------------------------------------------------------
\section{Implications for Audit Practice and Accountability}\label{sec:accountability}
Template-based audits are used as evidence that a system has been checked for
bias, including within internal audit processes that document a system's
development~\cite{raji2020}. Our results bear on what such evidence can
support.

\paragraph{Audit verdicts depend on choices that are rarely reported.} Across
our studies the verdict changed with the template type (at equal size, neutral
and emotional templates flag different biases, and race verdicts still differ
between them under a cluster-robust test), with the number of templates per
type, with how a multi-class output is reduced to a single score, and above all
with whether the significance test treats identity terms as fixed or as
sampled (for gender on the full EEC, 13 models flagged under the original test
and 2 under the robust one). Audit reports rarely state these choices, so a
reader cannot tell whether a verdict would survive a different, equally
reasonable design. The natural-sentence check adds a further reason for
caution: emotion-bearing templates produced large differences in uncertainty
that emotional content barely produced in real review sentences
(\S\ref{sec:natural}). The conditions under which a template audit observes a
model can therefore differ from those of the text the system will actually
process, and a verdict reached on templates need not carry over to that text.

\paragraph{Who relies on audit verdicts.} Different parties read the result of a
template audit for different purposes. Developers may use it to decide whether a
model needs further work; external auditors and regulators may treat it as
evidence that a system has been checked; and the people whose text a system
scores are affected by whether a bias is found and addressed, usually without
seeing how the audit was designed. A verdict that depends on unreported
template choices affects all of them, but not equally: a null result matters
most when it is used to justify deployment, and is least visible to those it
concerns. We do not study these parties directly; the point is that the
consequences of the measurement problem we document depend on who relies on
the result.

\paragraph{The risk of false reassurance.} Because identity effects are hard to
observe where a model is near-certain, an audit can return a null result partly
because its templates make the model confident. In our EEC replication, 4
gender and 5 race biases that were significant on the full EEC (paired $t$-test,
$\alpha=0.05/36$) were not significant on EEC-Neutral (2 and 3 under the
cluster-robust test), so an audit using only
those emotion-free templates, with that criterion, would not have detected them
(Table~\ref{tab:eec}). An equal-size comparison attributes two of these nine
misses to template type and the rest mainly to the small number of neutral
pairs (\S\ref{sec:eec}), so they show that the verdict depends on the
composition and size of the template set rather than isolating the role of
confidence. The three classifiers that are near-certain on
emotion-free text (FinancialBERT, finbert-tone and roberta-large-3cls)
illustrate the confidence mechanism more directly: on emotion-free templates, identity swaps changed
their output by almost nothing (mean $|\Delta_\ell|\leq 0.03$), whereas on
emotional templates the changes were considerably larger (0.08--0.90). For one
of them, finbert-tone, this included a significant gender bias that EEC-Neutral
did not detect; for the other two, neither template set yielded a significant
directional bias, so the difference concerns sensitivity rather than detected
bias (\S\ref{sec:construct}). These results hold for these models, templates and
statistical criterion. They show that a null result on a template set on which
the model is near-certain is uninformative, not that audits using emotion-free
templates generally give false reassurance. Such outcomes need not be
deliberate, since template sets are usually chosen without reference to the
audited model; but because template choice can affect whether a model passes an
audit, documenting the design is the safeguard.

\paragraph{Reporting recommendations.} In the spirit of documentation practices
such as model cards~\cite{mitchell2019}, we suggest that a template-based
counterfactual audit report the following.
\begin{enumerate}
  \item \emph{The model's uncertainty on each template set}, for example its
  median $1-\max_c P(c\mid x)$ and the share of templates on which the model is
  near-certain. A null result on templates where the model is near-certain is
  not evidence of fairness.
  \item \emph{Tests on a scale that does not saturate, alongside the
  probability scale.} A bias of known size is detected less often on the
  probability scale where models are confident (\S\ref{sec:injection}); the
  change in the model's logits avoids this, and if the verdicts on the two
  scales differ, both should be reported.
  \item \emph{Templates on which the model is uncertain}, in addition to standard
  ones, as a check on whether a null result is informative rather than as a
  better detector: if no template makes the model uncertain, a null result says
  little. Uncertainty-aware template selection with held-out person terms flagged
  a significant directional gender disparity in the most models in our study
  under the original paired test, but its advantage did not survive a
  cluster-robust test; we recommend it for checking informativeness, not for
  detection.
  \item \emph{Identity terms treated as a sampled factor.} Directional tests
  should account for the dependence among pairs that share templates and person
  terms (cluster-robust or mixed-effects tests), and report how many identity
  pairs the verdict rests on; with a handful of pairs, a verdict describes those
  terms rather than the group.
  \item \emph{Each template type, with enough templates, reported separately.}
  Four neutral templates are too few, the full EEC is dominated by emotional
  pairs, and bias directions can differ between types.
  \item \emph{The output measure and its justification} for the model's label
  scheme; valence-based scores can hide changes in multi-class models that a
  distributional measure reveals.
  \item \emph{Sensitivity, directional bias and label flips, separately.} Large
  sensitivity without a consistent direction, as in the binary models, is a
  different finding from a consistent bias.
  \item \emph{How identity categories are operationalised}, including which
  identities are left out and what the person terms are proxies for
  (\S\ref{sec:categories}).
  \item \emph{The design fixed before running, and all results reported.} We
  fixed our alternative LLM prompts, with recorded hashes, before running them;
  the same practice applies to template sets and output measures.
\end{enumerate}

% ----------------------------------------------------------------------------
\section{Limitations}
The association between uncertainty and identity sensitivity is correlational
and partly mechanical (\S\ref{sec:unc}); we have not manipulated uncertainty
directly, for example by temperature scaling or calibration. Identity
sensitivity measures how much the output changes, not whether the change is
harmful in a deployed setting. All templates are synthetic and English. The
natural-sentence check (\S\ref{sec:natural}) covers one domain (film reviews),
gender only, and an automatic swap with occasional errors; it supports the
within-model association but cannot confirm the between-condition contrast. In
languages with grammatical gender, substitutions require agreement changes that
our design avoids. The identity categories are binary for gender and use names
as proxies for race, with the limits discussed in \S\ref{sec:categories}. The neutrality of EEC-Neutral templates was not validated by human
annotators; we observe that several models do not treat them as neutral.
Uncertainty is measured as $1-\max_c P(c\mid x)$; normalised entropy and the
top-two margin give the same picture (\S\ref{sec:unc}), but other notions, such
as epistemic uncertainty estimated from ensembles, were not tested. The
predictive results concern probability-scale sensitivity; on a log-scale
measure, uncertainty did not predict multi-class models' direction beyond their
usual direction (\S\ref{sec:families}). A
cross-model control could not show that the gender gain of targeted sets
requires per-model selection (\S\ref{sec:targeted}). The targeted frames add mildly valenced clauses, so they
differ in content from the standard sets; the random-frame control holds the
clause pool fixed, but a detected bias reflects the model's behaviour in those
contexts, and its direction need not match that on other templates. The
targeted advantage in effect size is not uniform across comparisons, does not
clearly replicate in the four LLM variants, and the clause pool of 20 clauses
was written by us rather than sampled from natural text. The unseen template
families were also written by us and were tested on the 18 classifiers only; on
the logit-scale outcome, uncertainty did not add information beyond the label
scheme (\S\ref{sec:families}). The directional tests of Studies 1 and 3 rest on
few identity pairs (ten gender noun-phrase pairs and one name average in the
EEC, five gender and ten race pairs in Study~3); under a cluster-robust test
most gender verdicts and all Study~3 detections disappear, so these verdicts
describe the specific terms tested. The LLM results rest
on two models of one family and a reduced minimal-pair design. The direction
of the LLM effect was stable across three prompt wordings, but its size varied
several-fold, so LLM effect sizes should not be compared with those of the
classifiers; the EEC replication and Study~3 used the original wordings
only. The 18 models are public checkpoints rather than deployed
systems, and they are neither a random nor an independent sample of sentiment
models: they differ in architecture, label scheme and training domain, and
several share a base model. The filler and emotion clauses differ in wording
as well as in emotional content (\S\ref{sec:exp}), so Study~2 contrasts
matched clauses, not sentences identical except for emotion. Finally, the 20 EEC emotion words span a narrow range of affect
intensity, so our data cannot test whether the effect varies with the strength
of the emotion word; an exploratory analysis is included in the released code
but not interpreted here.

\section{Ethics Statement}
The study involves no human participants. Person terms serve as probes, not as
statements about individuals or groups, and the categories they stand for have
the limits set out in \S\ref{sec:categories}. A null result on a template set
should not be read as evidence that a model is fair (\S\ref{sec:accountability}).

\section{Data and Code Availability}
All code, the EEC-Neutral templates, the frozen LLM prompts and per-model results
are released (footnote~1). The Equity Evaluation Corpus is used under its
licence and is not redistributed; the release explains how to obtain it. The
natural sentences are not redistributed either; the released script rebuilds
them from the public IMDB corpus.

\section{Conclusion}
Counterfactual template audits are measurements, and their verdicts depend on
choices that audits rarely report. In a replication of the EEC, the template set
and, for multi-class models, the output measure determined whether a
directional bias was found and sometimes in which direction; the
neutral-sentence pattern reported for the original EEC systems is not a general
property of current sentiment classifiers. Where an audit can observe identity
effects is closely associated with predictive uncertainty: replacing an
emotion-bearing clause with a structurally matched filler clause increases
identity sensitivity in some models, has no detectable effect in others, and
decreases it in others, broadly tracking differences in uncertainty, and within
every system the same association holds on natural review sentences. In a test
fixed in advance on eight unseen template families, uncertainty measured with
held-out names predicted the direction of sensitivity on the primary,
probability-scale measure in 93\% of cells with a clear effect, including where
a model departed from its usual direction; on log-scale measures its value
beyond each model's usual direction was not established, so the prediction concerns where audits can see identity effects
rather than the models' underlying response. The association is partly mechanical and is
not evidence that uncertainty causes bias. Directional verdicts are fragile in
a further way: whether the significance test treats identity terms as fixed
or sampled changes the gender verdicts more than any template choice, and with
the EEC's few identity pairs most of them do not survive a cluster-robust test.
Uncertainty-aware template selection flags more gender disparities under the
original paired test, but not under the robust one. A small number of
significant flags should therefore not, on its own, be read as
reassurance, and neither should a large number be read as a finding about a
group rather than about the terms tested. Template-based audits should report model uncertainty and
the output measure, include uncertainty-aware templates alongside standard ones
as a check on whether a null result is informative, and cover each template type
with adequate power.

\subsubsection*{Acknowledgements and AI assistance.}
Generative AI tools were used to assist with writing and debugging the analysis
code and with language editing of the manuscript. The study design, analyses
and interpretation are the author's own; all code and results were checked by
the author, who takes full responsibility for the content.

% ----------------------------------------------------------------------------
\appendix
\section{EEC Results per Model}
\begin{table}[h]\centering\scriptsize
\caption{EEC replication per model (rows in the order of Table~\ref{tab:mp}).
Effect size $d_z$ of the paired difference on the logit scale (female minus
male; African-American minus European-American), on the full EEC, its
emotional templates, the original four neutral templates and EEC-Neutral;
$^{*}$ significant at $\alpha=0.05/36$ under the paired $t$-test of the original
protocol (cluster-robust counts in Table~\ref{tab:eec}). Ratio: mean
$|\Delta_\ell|$ on EEC-Neutral over emotional templates.}
\label{tab:eecfull}
\setlength{\tabcolsep}{2pt}
\begin{tabular}{l ccccc ccccc}
\toprule
& \multicolumn{5}{c}{Gender} & \multicolumn{5}{c}{Race} \\
\cmidrule(lr){2-6}\cmidrule(lr){7-11}
Model & Full & Emo. & Neu-4 & EEC-N & Ratio & Full & Emo. & Neu-4 & EEC-N & Ratio \\
\midrule
\inputrows{figures_tables/table_eec.tex}
\bottomrule
\end{tabular}
\end{table}

\clearpage
\section{Natural-Sentence Results per System}
\begin{table}[h]\centering\scriptsize
\caption{Natural-sentence check (\S\ref{sec:natural}): 1,000 IMDB sentences,
500 per affect stratum; classifiers in the order of Table~\ref{tab:mp}, then the
LLMs with the original prompts. $\rho$: within-system Spearman correlation
between uncertainty and identity sensitivity over sentences (all significant
after Holm correction). TV: mean identity sensitivity on low- and high-affect
sentences. Ratio: low over high affect, 95\% sentence-bootstrap CI (unadjusted);
$\downarrow$ significant, $\circ$ not significant after Holm correction across
the 22 systems. $u$ ratio: mean uncertainty, low over high affect. Bias:
directional gender bias on the logit scale ($\alpha=0.05/44$); F$>$M: female
versions scored more positively; --: none detected.}
\label{tab:natural}
\setlength{\tabcolsep}{3pt}
\begin{tabular}{l c c c l c c c c}
\toprule
System & $\rho$ & TV$_{\text{low}}$ & TV$_{\text{high}}$ & Ratio [95\% CI] & &
$u$ ratio & Bias$_{\text{low}}$ & Bias$_{\text{high}}$ \\
\midrule
\inputrows{figures_tables/table_natural.tex}
\bottomrule
\end{tabular}
\end{table}

\section{Exploratory Analyses of Study 4}\label{app:families}
Three checks added after the frozen analysis
qualify these results. First, restricting accuracy to clear cells does not
create the effect: using the point estimate in all 288 cells, the uncertainty
ratio gives the right direction in 81\% (72\% on the logit outcome). Second, the
two outcomes differ only in multi-class models. For binary models they are the
same measure and agree in every cell; in multi-class models, the logit outcome
points the same way, towards greater sensitivity with the emotion clause, in
93\% of clear cells, so a rule that assigns every multi-class model that
direction predicts it well. This matches the divergence between the
valence-based measure and TV in multi-class models reported in
\S\ref{sec:measure}, but we have not identified its cause: neutral-class
probability does not account for it (across models, $\rho=-0.34$ with the share
of cells in which the two outcomes disagree). Third, because TV is a
probability-scale measure, we repeated the test with one that does not
saturate, chosen before computing it: the mean absolute change in centred
log-probabilities, the distance between the two outputs in logit space. The
three binary models are unaffected (all 27 clear cells right). For the 15
multi-class models, the uncertainty ratio gives the right direction in 70\% of
clear cells (0.48--0.87), no better than each model's usual direction (73\%),
and on discordant cells it is at chance level (47\%, 0.14--0.85). For
multi-class models, uncertainty therefore predicts where a template change
makes identity effects more or less visible on the probability scale, but adds
nothing beyond the model's usual direction for the change in its logits. With
18 models and eight families as the independent units, all intervals here are
approximate.

\section{Unseen Template Families}
\begin{table}[h]\centering\scriptsize
\caption{The eight template families of the prospective test (\S\ref{sec:families}),
each with eight base sentences; the example shows the first base with the
emotion word \emph{furious} (or \emph{fury}) and its matched filler. Emotion
words per emotion: highest and lowest NRC intensity among the EEC words
(furious/annoyed, terrified/anxious, ecstatic/relieved, devastated/disappointed).}
\label{tab:families}
\setlength{\tabcolsep}{3pt}
\begin{tabular}{l p{0.3\linewidth} p{0.52\linewidth}}
\toprule
Family & Structure & Example (emotion / filler) \\
\midrule
\inputrows{figures_tables/table_families.tex}
\bottomrule
\end{tabular}
\end{table}

\section{LLM Minimal-Pair Results}
\begin{table}[h]\centering\scriptsize
\caption{Minimal-pair results for LLMs used as zero-shot classifiers (four
emotion words), for three prompt wordings per label scheme (p1: original; p2,
p3: alternatives fixed before running). Ratio: identity sensitivity, filler over
emotion, on the primary measure ($|\Delta_\ell|$ for binary, TV for three-class
prompts), 95\% two-way bootstrap CI; $\uparrow$ significant, $\circ$ not
significant after Holm correction across the 12 variants. Valence: the ratio on
$|\Delta_\ell|$. $\bar u$: mean uncertainty. $P_{\text{neu}}$: neutral
probability on filler sentences. Flips: share of pairs whose predicted label
changes, filler / emotion.}
\label{tab:llm}
\setlength{\tabcolsep}{3pt}
\begin{tabular}{l c c l c c cc c c}
\toprule
Model & Scheme & Wording & Ratio [95\% CI] & & Valence & $\bar u_{\text{fill}}$ &
$\bar u_{\text{emo}}$ & $P_{\text{neu}}$ & Flips \\
\midrule
\inputrows{figures_tables/table_llm.tex}
\bottomrule
\end{tabular}
\end{table}

\section{Model Identifiers}
\begin{table}[h]\centering\scriptsize
\caption{Hugging Face identifiers of the audited systems, with the short names
used in the paper and the number of output classes (for the LLMs, set by the prompt); classifiers in the order of
Table~\ref{tab:mp}.}
\label{tab:models}
\setlength{\tabcolsep}{3pt}
\begin{tabular}{l l c}
\toprule
Short name & Identifier & Classes \\
\midrule
\inputrows{figures_tables/table_models.tex}
\bottomrule
\end{tabular}
\end{table}

% ----------------------------------------------------------------------------
\begin{thebibliography}{99}
\bibitem{alnegheimish2022} Alnegheimish, S., Guo, A., Sun, Y.: Using Natural
Sentence Prompts for Understanding Biases in Language Models. In: Proc.\
NAACL-HLT, pp.\ 2824--2830 (2022)
\bibitem{bertrand2004} Bertrand, M., Mullainathan, S.: Are Emily and Greg More
Employable than Lakisha and Jamal? A Field Experiment on Labor Market
Discrimination. American Economic Review \textbf{94}(4), 991--1013 (2004)
\bibitem{blodgett2020} Blodgett, S.L., Barocas, S., Daum\'e III, H., Wallach,
H.: Language (Technology) is Power: A Critical Survey of ``Bias'' in NLP. In:
Proc.\ ACL (2020)
\bibitem{blodgett2021} Blodgett, S.L., Lopez, G., Olteanu, A., Sim, R.,
Wallach, H.: Stereotyping Norwegian Salmon: An Inventory of Pitfalls in
Fairness Benchmark Datasets. In: Proc.\ ACL-IJCNLP (2021)
\bibitem{caliskan2017} Caliskan, A., Bryson, J.J., Narayanan, A.: Semantics
derived automatically from language corpora contain human-like biases. Science
\textbf{356}(6334), 183--186 (2017)
\bibitem{card2020} Card, D., Henderson, P., Khandelwal, U., Jia, R., Mahowald, K., Jurafsky, D.:
With Little Power Comes Great Responsibility. In: Proc.\ EMNLP,
pp.\ 9263--9274 (2020)
\bibitem{czarnowska2021} Czarnowska, P., Vyas, Y., Shah, K.: Quantifying Social
Biases in NLP: A Generalization and Empirical Comparison of Extrinsic Fairness
Metrics. Transactions of the Association for Computational Linguistics
\textbf{9} (2021)
\bibitem{delobelle2022} Delobelle, P., Tokpo, E.K., Calders, T., Berendt, B.:
Measuring Fairness with Biased Rulers: A Comparative Study on Bias Metrics for
Pre-trained Language Models. In: Proc.\ NAACL-HLT, pp.\ 1693--1706 (2022)
\bibitem{dev2021} Dev, S., Monajatipoor, M., Ovalle, A., Subramonian, A.,
Phillips, J., Chang, K.-W.: Harms of Gender Exclusivity and Challenges in
Non-Binary Representation in Language Technologies. In: Proc.\ EMNLP,
pp.\ 1968--1994 (2021)
\bibitem{dev2022} Dev, S., Sheng, E., Zhao, J., Amstutz, A., Sun, J., Hou, Y.,
Sanseverino, M., Kim, J., Nishi, A., Peng, N., Chang, K.-W.: On Measures of
Biases and Harms in NLP. In: Findings of the ACL: AACL-IJCNLP 2022,
pp.\ 246--267 (2022)
\bibitem{dixon2018} Dixon, L., Li, J., Sorensen, J., Thain, N., Vasserman, L.:
Measuring and Mitigating Unintended Bias in Text Classification. In: Proc.\
AAAI/ACM Conference on AI, Ethics, and Society (AIES) (2018)
\bibitem{dror2018} Dror, R., Baumer, G., Shlomov, S., Reichart, R.: The Hitchhiker's Guide to
Testing Statistical Significance in Natural Language Processing. In: Proc.\
56th Annual Meeting of the ACL, pp.\ 1383--1392 (2018)
\bibitem{field2021} Field, A., Blodgett, S.L., Waseem, Z., Tsvetkov, Y.: A
Survey of Race, Racism, and Anti-Racism in NLP. In: Proc.\ ACL-IJCNLP,
pp.\ 1905--1925 (2021)
\bibitem{gaddis2017} Gaddis, S.M.: How Black Are Lakisha and Jamal? Racial
Perceptions from Names Used in Correspondence Audit Studies. Sociological
Science \textbf{4}, 469--489 (2017)
\bibitem{gallegos2024} Gallegos, I.O., Rossi, R.A., Barrow, J., Tanjim, M.M., Kim, S., Dernoncourt,
F., Yu, T., Zhang, R., Ahmed, N.K.: Bias and Fairness in Large Language
Models: A Survey. Computational Linguistics \textbf{50}(3), 1097--1179 (2024)
\bibitem{garg2019} Garg, S., Perot, V., Limtiaco, N., Taly, A., Chi, E.H.,
Beutel, A.: Counterfactual Fairness in Text Classification through Robustness.
In: Proc.\ AAAI/ACM Conference on AI, Ethics, and Society (AIES) (2019)
\bibitem{goldfarb2021} Goldfarb-Tarrant, S., Marchant, R., Mu\~noz S\'anchez,
R., Pandya, M., Lopez, A.: Intrinsic Bias Metrics Do Not Correlate with
Application Bias. In: Proc.\ ACL-IJCNLP (2021)
\bibitem{goldfarb2023} Goldfarb-Tarrant, S., Ungless, E., Balkir, E., Blodgett, S.L.: This Prompt
is Measuring \textless MASK\textgreater: Evaluating Bias Evaluation in Language Models. In:
Findings of the ACL: ACL 2023, pp.\ 2209--2225 (2023)
\bibitem{guo2017} Guo, C., Pleiss, G., Sun, Y., Weinberger, K.Q.: On
Calibration of Modern Neural Networks. In: Proc.\ ICML (2017)
\bibitem{huang2020} Huang, P.-S., Zhang, H., Jiang, R., Stanforth, R., Welbl,
J., Rae, J., Maini, V., Yogatama, D., Kohli, P.: Reducing Sentiment Bias in
Language Models via Counterfactual Evaluation. In: Findings of EMNLP (2020)
\bibitem{jacobs2021} Jacobs, A.Z., Wallach, H.: Measurement and Fairness. In:
Proc.\ ACM Conference on Fairness, Accountability, and Transparency (FAccT),
pp.\ 375--385 (2021)
\bibitem{kiritchenko2018} Kiritchenko, S., Mohammad, S.M.: Examining Gender and
Race Bias in Two Hundred Sentiment Analysis Systems. In: Proc.\ 7th Joint
Conference on Lexical and Computational Semantics (*SEM), pp.\ 43--53 (2018)
\bibitem{krishna2022} Krishna, S., Gupta, R., Verma, A., Dhamala, J., Pruksachatkun, Y.,
Chang, K.-W.: Measuring Fairness of Text Classifiers via Prediction Sensitivity.
In: Proc.\ 60th Annual Meeting of the ACL (Volume 1: Long Papers),
pp.\ 5830--5842 (2022)
\bibitem{kusner2017} Kusner, M.J., Loftus, J., Russell, C., Silva, R.:
Counterfactual Fairness. In: Advances in Neural Information Processing Systems
(NeurIPS) (2017)
\bibitem{kuzucu2024} Kuzucu, S., Cheong, J., Gunes, H., Kalkan, S.: Uncertainty as a Fairness
Measure. Journal of Artificial Intelligence Research \textbf{81}, 307--335 (2024)
\bibitem{larson2017} Larson, B.: Gender as a Variable in Natural-Language
Processing: Ethical Considerations. In: Proc.\ First ACL Workshop on Ethics
in Natural Language Processing (2017)
\bibitem{maas2011} Maas, A.L., Daly, R.E., Pham, P.T., Huang, D., Ng, A.Y.,
Potts, C.: Learning Word Vectors for Sentiment Analysis. In: Proc.\ 49th Annual
Meeting of the ACL: Human Language Technologies, pp.\ 142--150 (2011)
\bibitem{mitchell2019} Mitchell, M., Wu, S., Zaldivar, A., Barnes, P.,
Vasserman, L., Hutchinson, B., Spitzer, E., Raji, I.D., Gebru, T.: Model Cards
for Model Reporting. In: Proc.\ Conference on Fairness, Accountability, and
Transparency (FAT*) (2019)
\bibitem{mohammad2018} Mohammad, S.M.: Word Affect Intensities. In: Proc.\ 11th
International Conference on Language Resources and Evaluation (LREC) (2018)
\bibitem{orgad2022} Orgad, H., Belinkov, Y.: Choose Your Lenses: Flaws in Gender Bias
Evaluation. In: Proc.\ 4th Workshop on Gender Bias in Natural Language
Processing (GeBNLP), pp.\ 151--167 (2022)
\bibitem{pleiss2017} Pleiss, G., Raghavan, M., Wu, F., Kleinberg, J., Weinberger, K.Q.: On Fairness
and Calibration. In: Advances in Neural Information Processing Systems 30,
pp.\ 5680--5689 (2017)
\bibitem{prabhakaran2019} Prabhakaran, V., Hutchinson, B., Mitchell, M.:
Perturbation Sensitivity Analysis to Detect Unintended Model Biases. In: Proc.\
EMNLP-IJCNLP (2019)
\bibitem{raji2020} Raji, I.D., Smart, A., White, R.N., Mitchell, M., Gebru,
T., Hutchinson, B., Smith-Loud, J., Theron, D., Barnes, P.: Closing the AI
Accountability Gap: Defining an End-to-End Framework for Internal Algorithmic
Auditing. In: Proc.\ Conference on Fairness, Accountability, and Transparency
(FAT*) (2020)
\bibitem{ribeiro2020} Ribeiro, M.T., Wu, T., Guestrin, C., Singh, S.: Beyond
Accuracy: Behavioral Testing of NLP Models with CheckList. In: Proc.\ ACL
(2020)
\bibitem{selvam2023} Selvam, N.R., Dev, S., Khashabi, D., Khot, T., Chang,
K.-W.: The Tail Wagging the Dog: Dataset Construction Biases of Social Bias
Benchmarks. In: Proc.\ ACL (Short Papers), pp.\ 1373--1386 (2023)
\bibitem{seshadri2022} Seshadri, P., Pezeshkpour, P., Singh, S.: Quantifying
Social Biases Using Templates is Unreliable. In: NeurIPS Workshop on Trustworthy
and Socially Responsible Machine Learning (2022). arXiv:2210.04337
\bibitem{smith2022} Smith, E.M., Hall, M., Kambadur, M., Presani, E., Williams, A.: ``I'm sorry to
hear that'': Finding New Biases in Language Models with a Holistic Descriptor
Dataset. In: Proc.\ EMNLP, pp.\ 9180--9211 (2022)
\bibitem{vanderwal2024} van der Wal, O., Bachmann, D., Leidinger, A., van Maanen, L., Zuidema, W.,
Schulz, K.: Undesirable Biases in NLP: Addressing Challenges of Measurement.
Journal of Artificial Intelligence Research \textbf{79}, 1--40 (2024)
\bibitem{zhang2021} Zhang, C., Zhao, J., Zhang, H., Chang, K.-W., Hsieh, C.-J.:
Double Perturbation: On the Robustness of Robustness and Counterfactual Bias
Evaluation. In: Proc.\ NAACL-HLT, pp.\ 3899--3916 (2021)
\end{thebibliography}

\end{document}
